\documentclass[conference]{IEEEtran}

\usepackage{amsmath}
\usepackage{cite}

\title{Data Re-Uploading and the Fourier Spectrum of Quantum Neural Networks}

\author{\IEEEauthorblockN{Rana Muhammad Hamza}
\IEEEauthorblockA{Roll Number: 25280103\\
QML Course\\
Instructor: Dr.\ Faryad}}

\begin{document}
\maketitle

\begin{abstract}
This project will test how data re-uploading and circuit depth shape the
Fourier components representable by quantum neural networks (QNNs). We will
fit controlled one- and two-dimensional periodic targets with one to six
encoding layers, compare serial and parallel feature encodings, and evaluate
a two-qubit classifier on synthetic spirals. Fourier-feature linear
regression and a parameter-matched multilayer perceptron (MLP) will provide
classical baselines. Noiseless simulations will establish spectra and
generalization errors; a trained one-qubit regressor will additionally be
evaluated on IBM quantum hardware, subject to account and backend access.
The analysis will test predicted frequency support while distinguishing
theoretically allowed harmonics from modes actually present after training.
\end{abstract}

\begin{IEEEkeywords}
quantum machine learning, data re-uploading, Fourier spectrum, variational
quantum circuits, quantum neural networks
\end{IEEEkeywords}

\section{Motivation and Research Question}
Data re-uploading alternates input-dependent quantum gates with trainable
operations, allowing a compact circuit to build richer functions through
repeated encoding \cite{perez,schuld}. The number of encoding repetitions
therefore provides a principled control over the candidate frequencies of a
QNN. This project asks: \emph{How do encoding repetitions and layer count
determine the frequencies a QNN can represent, and how do serial and parallel
feature encodings affect that representation and its resource cost?}

For a feature encoded by $L$ occurrences of rotations
$R_P(x)=\exp(-ixP/2)$, where $P$ is a Pauli operator, the expected output has
candidate integer harmonics $k\in\{-L,\ldots,L\}$. In multiple dimensions,
the support is a set of frequency vectors; the circuit structure can restrict
that set, and optimization can suppress otherwise allowed coefficients.
Accordingly, the study will report both candidate support and measured
Fourier coefficients rather than equating the two.

\section{Data and Experimental Design}
\subsection{Synthetic Regression}
Inputs lie on the periodic domain $[-\pi,\pi)$ (and its two-dimensional
product). The target functions are
\begin{align}
f_1(x) &= 0.5\cos x+0.3\sin(2x)+0.2\cos(3x),\\
f_2(x_1,x_2) &= 0.4\sin x_1+0.3\cos(2x_2)
                  +0.3\sin(x_1-x_2).
\end{align}
Both functions are bounded by one and have known, nontrivial spectra. Models
with one through six re-uploading layers will be trained on fixed,
seeded uniform samples (128 one-dimensional and 2,048 two-dimensional
points) and evaluated on separate dense uniform grids. Each configuration
will use data-generation seed 42, five initialization seeds (0--4), and
full-batch Adam (learning rate 0.01, at most 1,000 updates) to minimize MSE.
The one-dimensional grid will contain 256 points and the two-dimensional
grid will contain $64\times64$ points. Noiseless predictions on these grids
will be transformed with a discrete Fourier transform (DFT); grid
refinement will check suspected out-of-support modes for aliasing.

\subsection{Quantum Models and Encoding Ablation}
The primary models are variational circuits with Pauli-rotation data
encoding, trainable single-qubit rotations, and entangling operations where
applicable. Layer count $L\in\{1,\ldots,6\}$ is the controlled depth variable.
For one feature, a one-qubit circuit re-uploads $x$ once per layer. For two
features, the \emph{serial} circuit applies both feature rotations in sequence
on one qubit per layer; the \emph{parallel} circuit encodes one feature on
each of two qubits per layer and entangles them. The expectation of a Pauli-$Z$
observable supplies the regression output or classification score. Encoding
comparisons will match trainable-parameter counts as closely as circuit
structure permits and will always report qubits, depth, two-qubit gates, and
trainable parameters. For both two-feature layouts, the candidate support
will be tested against coordinate-wise bounds $|k_j|\leq L$ and any
circuit-specific restrictions.

\subsection{Classification and Baselines}
The classification task uses a fixed-seed, balanced two-class spiral dataset
with 1,000 examples per class. Five-fold stratified cross-validation will fit
any scaling or preprocessing on training folds only. Regression baselines
are Fourier-feature linear regression, using harmonics up to the QNN's
candidate cutoff, and an MLP; for spirals, the same Fourier features feed
logistic regression and the MLP is a classifier. MLP parameter counts will
match the QNN as closely as possible. Baseline feature maps and tuning
procedures will be documented. Report mean and standard deviation of
held-out accuracy and macro F1 for classification and MSE for regression.

\section{Hypothesis, Falsifier, and Analysis}
The primary hypothesis is that repeated Pauli encodings constrain candidate
Fourier frequencies according to the number and placement of encoding
occurrences, while optimization and circuit connectivity determine which
candidate modes acquire appreciable coefficients. For each trained model,
the DFT will report coefficient magnitudes and locations alongside held-out
MSE. The serial/parallel ablation will compare spectra and predictive
performance at matched layer and parameter budgets, with resource counts
reported separately rather than conflated.

The falsifier is a reproducible Fourier coefficient outside the
architecture's predicted frequency support in a noiseless simulation. A
coefficient will count as evidence only if its magnitude exceeds
$\max(10^{-6},\,0.01\,c_{\max})$, where $c_{\max}$ is the largest measured
coefficient, and it persists under a refined evaluation grid and independent
parameter initializations. Before drawing that conclusion, we will verify
the encoding-angle convention, frequency normalization, numerical precision,
and absence of sampling aliasing. Noise and finite-shot experiments will be
reported separately and will not be used alone to claim a violation of the
noiseless prediction.

\section{Validation and Hardware Evaluation}
We will use the course NB11 capstone harness for fixed folds, leakage-free
preprocessing, baselines, ablations, and noise sweeps. Seeds, generated data,
model configurations, package versions, and results will be saved with the
code. The hardware study will execute the trained three-layer, one-qubit periodic
regressor on an IBM backend at 64 evenly spaced inputs with 1,024 shots per
input, compare its finite-shot curve with the noiseless simulator curve, and
record backend, calibration context, and circuit resources. Access
limitations will be reported explicitly rather than substituted with
simulated data.

\section{Week-by-Week Plan and Submissions}
\begin{table}[!t]
\caption{Project milestones and planned evidence}
\label{tab:timeline}
\centering
\footnotesize
\begin{tabular}{p{0.20\columnwidth}p{0.72\columnwidth}}
\hline
\textbf{Date} & \textbf{Plan and submission}\\
\hline
Nov.\ 2 (0\%) & Confirm project/group; enter names and GitHub URL on sign-up tab.\\
Nov.\ 8 (5\%) & Submit 1--2 page synopsis: question, falsifier, data, models, baselines, weekly plan.\\
Nov.\ 15 (15\%) & Show data pipeline, fixed folds, classical baselines, first QNN; submit introduction and methods draft and GitHub code.\\
Nov.\ 22 (25\%) & Show simulator results, noise model, and ablations; submit results draft and code.\\
Nov.\ 29 (35\%) & Present for first-half groups; have IBM hardware result and final talk ready.\\
Dec.\ 6 (20\%) & Remaining presentations; submit final LaTeX report (PDF and source) and tagged GitHub release.\\
\hline
\end{tabular}
\end{table}
Every Sunday, the group will present its progress, current report draft, and
GitHub code. The final report will follow the two-column IEEEtran or
REV\TeX{} 4.2 format and include resource counts and author contributions;
the presentation target is 12 minutes plus 5 minutes for questions.

\begin{thebibliography}{2}
\bibitem{perez}
A. P\'erez-Salinas, A. Cervera-Lierta, E. Gil-Fuster, and J. I. Latorre,
``Data re-uploading for a universal quantum classifier,''
\emph{Quantum}, vol. 4, p. 226, 2020, doi: 10.22331/q-2020-02-06-226.

\bibitem{schuld}
M. Schuld, R. Sweke, and J. J. Meyer, ``Effect of data encoding on the
expressive power of variational quantum-machine-learning models,''
\emph{Phys. Rev. A}, vol. 103, p. 032430, 2021,
doi: 10.1103/PhysRevA.103.032430.
\end{thebibliography}

\end{document}
